% ============================================================
%  CAPÍTULO 1 — PRESENTACIÓN DE LA EMPRESA (Subdoc. 1)
%  Contenido: empresa, trayectoria, innovación, equipo
% ============================================================
%  FUENTES:
%  - Subdoc. 1 (Bases_Administrativas.md §1532-1537):
%      §1534 reseña/trayectoria/capacidades/lineas de negocio.
%      §1535 estructura org., dotación, certificaciones y alianzas.
%      §1536 experiencia en la industria del caso y complejidad equivalente.
%      §1537 gobierno interno de calidad, seguridad y gestión del conocimiento.
%  - T-22, Informe 1 (§1838): presentación de la empresa.
%  - Subdoc. 12 (§1627): estructura del equipo / gobernanza (ref. transversal).
%  Ver FUENTES.md para el detalle completo.
\chapter{Presentación de la Empresa (Subdoc.\ 1)}

\nota{Este capítulo corresponde al subdoc.\ 1 del Formulario T-22. Aquí se presenta la empresa proponente: quiénes son, qué hacen, qué experiencia tienen y por qué están en condiciones de responder a esta licitación. No se describe el problema del caso ni la solución; eso va en los capítulos 2 y 3. El evaluador busca ver que la empresa es real, que tiene trayectoria relevante y que cuenta con el equipo adecuado.}

% -----------------------------------------------------------
\section{Identificación y Perfil Corporativo}
% -----------------------------------------------------------

\nota{Datos básicos de la empresa: razón social, RUT, giro, dirección, representante legal, fecha de constitución, tamaño según SERVEL (micro/pequeña/mediana/grande). Si la empresa es nueva o no tiene licitaciones previas, explicitarlo sin esconderlo. Si hay wording del caso que describes como empresa existente, ajustar la narrativa a la empresa real.}

\pendiente

% -----------------------------------------------------------
\section{Filosofía Estratégica: Misión, Visión y Sostenibilidad}
% -----------------------------------------------------------

\nota{Redactar misión, visión y compromiso ESG (ambiental, social, gobernanza). Si la empresa no tiene declaraciones formales, formularlas de forma creíble alineadas al proyecto. El evaluador busca ver coherencia entre lo que la empresa dice ser y lo que propone hacer.}

\subsection{Misión}
\pendiente

\subsection{Visión}
\pendiente

\subsection{Compromiso de Sostenibilidad y Gobernanza (ESG)}
\pendiente

% -----------------------------------------------------------
\section{Trayectoria y Casos de Éxito}
% -----------------------------------------------------------

\nota{Describir proyectos anteriores relevantes a esta licitación. Idealmente 2 a 4 casos. Para cada uno: nombre del cliente, alcance del proyecto, tecnologías usadas, resultados medibles (porcentajes, tiempos, ahorros). Si no hay casos formales, usar proyectos académicos o personales, pero ser transparente. Lo que el evaluador evalúa aquí es madurez técnica demostrada, no antigüedad de la empresa.}

\subsection{Caso de éxito 1 — Trazabilidad sanitaria}
\pendiente
\nota{Describir un caso donde la empresa haya implementado o participado en un sistema de trazabilidad de productos, idealmente con componentes de cadena de frío o cumplimiento sanitario. Si no existe, plantear un proyecto académico o de investigación relevante.}

\subsection{Caso de éxito 2 — Logística / última milla}
\pendiente
\nota{Describir un caso donde la empresa haya trabajado en optimización de rutas, gestión de flota, o distribución de consumo masivo. Puede ser un proyecto académico, una consultoría, o una contribución a software open-source relevante.}

% -----------------------------------------------------------
\section{Credenciales Organizacionales y Madurez TI}
% -----------------------------------------------------------

\nota{Certificaciones (ISO 9001, ISO 27001, SOC 2, si existen), stack tecnologico dominado por la empresa, frameworks de gestion de proyectos (Scrum, PMI), y cualquier otra credencial que respalde la capacidad tecnica. Si no hay certificaciones formales, describir la madurez tecnica del equipo: lenguajes, frameworks, cloud, bases de datos, herramientas DevOps.}

\pendiente

% -----------------------------------------------------------
\section{Capacidad Instalada para Proyectos Complejos}
% -----------------------------------------------------------

\nota{Describir la capacidad operativa: cuántos desarrolladores, infraestructura disponible, experiencia con proyectos desimilar envergadura al de esta licitación. Si la empresa es unipersonal o pequeño equipo, ser honesto y enfatizar la calidad del trabajo sobre el tamaño.}

\pendiente

% -----------------------------------------------------------
\section{Estructura del Equipo de Trabajo y Gobernanza}
% -----------------------------------------------------------

\nota{Presentar el equipo que trabajará en la propuesta y su ejecución. Nombrar roles (líder de proyecto, analista, desarrollador, tester, etc.), experiencia de cada uno y cómo se coordinarán. El evaluador busca ver que hay personas suficientes y capacitadas para el alcance propuesto. Incluir una tabla con roles, nombre y responsabilidad principal.}

\begin{tabularx}{\textwidth}{|L{3.8cm}|L{3.2cm}|X|}
\hline
\rowcolor{lafroxClaro}
\textbf{Rol} & \textbf{Nombre} & \textbf{Responsabilidad principal} \\
\hline
Jefe de Proyecto & Alex Aravena & Dirección técnica, coordinación, interfaz con el cliente \\
\hline
Arquitecto de Solución & Bastián Trejo & Diseño de arquitectura, selección tecnológica \\
\hline
Encargado de Seguridad de la Información & Álvaro Catalán & Ciberseguridad, trazabilidad de datos, cumplimiento \\
\hline
Líder de Datos & Leandro Chamorro & Modelo de datos, migración, gobernanza y analítica \\
\hline
Líder de Desarrollo & \pendiente & Implementación de componentes \\
\hline
Líder de Calidad & \pendiente & Planificación y ejecución de pruebas \\
\hline
Líder de Operación / SRE & \pendiente & Disponibilidad, observabilidad, operación \\
\hline
Líder de Implantación y Gestión del Cambio & \pendiente & Despliegue, capacitación, adopción en terreno \\
\hline
\end{tabularx}
